% Part: first-order-logic
% Chapter: introduction
% Section: syntax

\documentclass[../../../include/open-logic-section]{subfiles}

\begin{document}

\olfileid{fol}{int}{syn}

\olsection{نحو}

لومړے بايد دا په دقيق ډول وټاکو چې د سمبولونو کوم تارونه د لومړۍ
درجې منطق !!{sentence}s ګڼل کېږي۔ دا کار به وروسته وکړو؛ اوس به
يوازې د بېلګو له مخې مخکښې ځو۔ د لومړۍ درجې منطق بنسټيزې جوړښتي
خښتې---يعنې لغتځونډ---په دوو برخو وېشل کېږي۔ لومړۍ برخه هغه
سمبولونه دي چې په وسيله ئې ځانګړي څيزونه ټاکو يا د هغوئ په باب
ځانګړې خبرې کوو۔ څيزونه په !!{constant}s ټاکو، او بيا د ټاکلو شوو
څيزونو په باب په !!{predicate}s څه وايو۔ مثلاً ښايي~$\Obj a$ د
!!a{constant} په توګه د يوه څيز د ټاکلو دپاره وکاروو، او بيا د
!!{sentence}~$\Atom{\Obj P}{\Obj a}$ په وسيله د هغه په باب څه
ووايو۔ که د ``$\Obj a$'' او ``$\Obj P$'' معناوې په ذهن کښې ولرئ،
$\Atom{\Obj P}{\Obj a}$ د انګرېزي ژبې د يوې جملې په توګه لوستلے
شئ (او غالباً مو د صوري منطق د لومړۍ زده کړې پر وخت همداسې کړي
دي)۔ کله چې د لومړۍ درجې منطق داسې ساده !!{sentence}s ولرو، د
لغتځونډ په دويمه برخه، يعنې منطقي سمبولونو (نښلوونکو او کميت
ټاکونکو)، ترې نورې پېچلې جملې جوړولے شو۔ نو مثلاً داسې عبارتونه
جوړولے شو لکه $(\Atom{\Obj P}{\Obj a} \land \Atom{\Obj Q}{\Obj b})$
يا~$\lexists[\Obj x][\Atom{\Obj P}{\Obj x}]$۔

د معنٰی پوهنې د دقيقو تعريفونو او د هغو د !!{derivation} د نظامونو
د قاعدو د ورکولو دپاره چې سخته مابعدمنطقي څېړنه ئې غواړي، تر هر څه
مخکښې بايد په دقيق ډول تعريف کړو چې د لومړۍ درجې منطق
!!a{sentence} څه ګڼل کېږي۔ بنسټيزه مفکوره دومره اسانه ده: ځينې ساده
!!{sentence}s يوازې له !!{predicate}s او !!{constant}s څخه جوړولے
شو، لکه~$\Atom{\Obj P}{\Obj a}$۔ بيا له دغو څخه د نښلوونکو او کميت
ټاکونکو په وسيله پېچلې جملې جوړوو۔ خو هغه قواعد په دقيق ډول کوم دي
چې له مخې ئې د پېچلو !!{sentence}s جوړول راته جايز دي؟ دا بايد
وټاکل شي، کنه ``د لومړۍ درجې منطق !!{sentence}'' مو په کافي دقت نۀ
دے تعريف کړے۔ څو مسئلې شته۔ لومړۍ دا چې سم تارونه د !!{sentence}s
په توګه وشمېرو۔ دويمه دا چې کار په داسې ډول وکړو چې د
\emph{ټولو} !!{sentence}s په باب رياضيکي ثبوتونه ورکولے شو۔ په پاے
کښې بايد پر !!{sentence}s د ځينو لومړنيو عملياتو دقيق تعريفونه هم
ورکړو، لکه ``په~$!A$ کښې هر~$\Obj x$ په~$\Obj b$ بدل کړه''۔ ستونزه
دا ده چې د لومړۍ درجې منطق کميت ټاکونکي او !!{variable}s دا کار
بشپړ څرګند نۀ پرېږدي۔ مثلاً ايا
$\lexists[\Obj x][\Atom{\Obj P}{\Obj a}]$ بايد !!a{sentence} وګڼل
شي؟ او $\lexists[\Obj x][\lexists[\Obj x][\Atom{\Obj P}{\Obj x}]]$
څنګه؟ ``په $(\Atom{\Obj P}{\Obj x} \land \lexists[\Obj x][\Atom{\Obj P}{\Obj x}])$
کښې $\Obj x$ په~$\Obj b$ بدل کړه'' پايله بايد څه وي؟

\end{document}
